% problem_id: S001-lemma-recover-tail-graded-module-general
% source_id: S001 (Stacks Project)
% source_file: coherent.tex
% statement_locator: coherent.tex lines 4080--4097
% proof_locator: coherent.tex lines 4099--4197
% env: lemma
% pairing: tight (verified)
% status: complete (statement + author proof verbatim + reason + locators)
%
\begin{lemma}
\label{lemma-recover-tail-graded-module-general}
Let $A$ be a Noetherian graded ring and let $d$ be the lcm of generators
of $A$ over $A_0$. Let $M$ be a finite graded $A$-module.
Set $X = \text{Proj}(A)$ and let $\widetilde{M}$ be
the quasi-coherent $\mathcal{O}_X$-module on $X$ associated to $M$.
Let $k \in \mathbf{Z}$.
\begin{enumerate}
\item $N' = \bigoplus_{n \geq k} H^0(X, \widetilde{M(n)})$
is a finite $A$-module,
\item $N = \bigoplus_{n \geq k} H^0(X, \widetilde{M}(n))$
is a finite $A$-module,
\item there is a canonical map $N \to N'$,
\item if $k$ is small enough there is a canonical map $M \to N'$,
\item the map $M_n \to N'_n$ is an isomorphism for $n \gg 0$,
\item $N_n \to N'_n$ is an isomorphism for $d | n$.
\end{enumerate}
\end{lemma}

\begin{proof}
The map $N \to N'$ in (3) comes from
Constructions, Equation (\ref{constructions-equation-multiply-more-generally})
by taking global sections.

\medskip\noindent
By
Constructions, Equation
(\ref{constructions-equation-global-sections-more-generally})
there is a map of graded $A$-modules
$M \to \bigoplus_{n \in \mathbf{Z}} H^0(X, \widetilde{M(n)})$.
If the generators of $M$ sit in degrees $\geq k$, then the image
is contained in the submodule
$N' \subset \bigoplus_{n \in \mathbf{Z}} H^0(X, \widetilde{M(n)})$
and we get the map in (4).

\medskip\noindent
By Algebra, Lemmas \ref{algebra-lemma-graded-Noetherian} and
\ref{algebra-lemma-S-plus-generated} the ring $A_0$ is Noetherian
and $A$ is generated over $A_0$ by finitely many elements
$f_1, \ldots, f_r$ homogeneous of positive degree.
Let $d = \text{lcm}(\deg(f_i))$. Then we see that (6) holds
for example by
Constructions, Lemma \ref{constructions-lemma-where-invertible}.

\medskip\noindent
Because $M$ is a finite $A$-module we see that
$\widetilde{M}$ is a finite type $\mathcal{O}_X$-module,
i.e., a coherent $\mathcal{O}_X$-module.
Thus part (2) follows from Lemma \ref{lemma-coherent-on-proj-general}.

\medskip\noindent
We will deduce (1) from (2) using a trick. For $q \in \{0, \ldots, d - 1\}$
write
$$
{}^qN = \bigoplus\nolimits_{n + q \geq k} H^0(X, \widetilde{M(q)}(n))
$$
By part (2) these are finite $A$-modules. The Noetherian ring $A$
is finite over $A^{(d)} = \bigoplus_{n \geq 0} A_{dn}$, because
it is generated by $f_i$ over $A^{(d)}$ and $f_i^d \in A^{(d)}$.
Hence ${}^qN$ is a finite $A^{(d)}$-module.
Moreover, $A^{(d)}$ is Noetherian (follows from
Algebra, Lemma \ref{algebra-lemma-dehomogenize-finite-type}).
It follows that the $A^{(d)}$-submodule
${}^qN^{(d)} = \bigoplus_{n \in \mathbf{Z}} {}^qN_{dn}$
is a finite module over $A^{(d)}$. Using the isomorphisms
$\widetilde{M(dn + q)} = \widetilde{M(q)}(dn)$ we can write
$$
N' =
\bigoplus\nolimits_{q \in \{0, \ldots, d - 1\}}
\bigoplus\nolimits_{dn + q \geq k}
H^0(X, \widetilde{M(q)}(dn)) =
\bigoplus\nolimits_{q \in \{0, \ldots, d - 1\}} {}^qN^{(d)}
$$
Thus $N'$ is finite over $A^{(d)}$ and a fortiori finite over $A$.
Thus (1) is true.

\medskip\noindent
Let $K$ be a finite $A$-module such that $\widetilde{K} = 0$. We claim
that $K_n = 0$ for $d|n$ and $n \gg 0$. Arguing as above we see that
$K^{(d)}$ is a finite $A^{(d)}$-module. Let
$x_1, \ldots, x_m \in K$ be homogeneous generators of $K^{(d)}$
over $A^{(d)}$, say sitting in degrees $d_1, \ldots, d_m$ with $d | d_j$.
For each $i$ and $j$ there exists an $n_{ij} \geq 0$ such that
$f_i^{n_{ij}} x_j = 0$ in $K_{d_j + n_{ij}}$: if not then
$x_j/f_i^{d_i/\deg(f_i)} \in K_{(f_i)}$ would not be zero, i.e.,
$\widetilde{K}$ would not be zero. Here we use that $\deg(f_i) | d | d_j$
for all $i, j$. We conclude that $K_n$ is zero for
$n$ with $d | n$ and $n > \max_j (d_j + \sum_i n_{ij} \deg(f_i))$
as every element of $K_n$ is a sum of terms where each term is a
monomials in the $f_i$ times one of the $x_j$ of total degree $n$.

\medskip\noindent
To finish the proof, we have to show that $M \to N'$ is an isomorphism
in all sufficiently large degrees. The map $N \to N'$ induces an
isomorphism $\widetilde{N} \to \widetilde{N'}$ because on the affine
opens $D_+(f_i) = D_+(f_i^d)$ the corresponding modules are isomorphic:
$N_{(f_i)} \cong N_{(f_i^d)} \cong N'_{(f_i^d)} \cong N'_{(f_i)}$
by property (6).
By Properties, Lemma \ref{properties-lemma-proj-quasi-coherent}
we have a canonical isomorphism $\widetilde{N} \to \widetilde{M}$.
The composition $\widetilde{N} \to \widetilde{M} \to \widetilde{N'}$
is the isomorphism above (proof omitted; hint: look on standard
affine opens to check this). Thus the map $M \to N'$ induces an isomorphism
$\widetilde{M} \to \widetilde{N'}$.
Let $K = \Ker(M \to N')$ and $Q = \Coker(M \to N')$.
Recall that the functor
$M \mapsto \widetilde{M}$ is exact, see
Constructions, Lemma \ref{constructions-lemma-proj-sheaves}.
Hence we see that $\widetilde{K} = 0$ and $\widetilde{Q} = 0$.
By the result of the previous paragraph we see that $K_n = 0$ and
$Q_n = 0$ for $d | n$ and $n \gg 0$. At this point we finally see
the advantage of using $N'$ over $N$: the functor $M \leadsto N'$
is compatible with shifts (immediate from the construction).
Thus, repeating the whole argument with $M$ replaced by $M(q)$
we find that $K_n = 0$ and $Q_n = 0$ for $n \equiv q \bmod d$
and $n \gg 0$. Since there are
only finitely many congruence classes modulo $n$ the proof is finished.
\end{proof}
